\sectionstart
\section{版本演进与后续工作}
\subsection{3.2 系列的工程收敛}
3.2 将已有转账、账单与订阅能力连接到可解释财务体检，并补充术语解释、长辈友好显示、模型中心与注册账户。后续迭代集中在三个方面：认证交互与输入一致性，官方跑腿组件的授权、部署与错误处理，以及模型自由对话、来源和超时语义。3.2.9 的主要落点是千问普通对话与连接测试的稳定请求策略。

版本号并不增加已经验证的能力范围。水墨视觉更新改善项目表达；财务算法、交易状态与外部服务仍按本文定义运行。发布前应把前端、后端、配置、数据库兼容与验证证据绑定到同一可追溯快照。

\subsection{优先完成的工程改进}
\begin{enumerate}
\item \textbf{统一动作确认契约。} 将普通转账、多步骤任务、订阅、卡片与理财变更逐一映射到相同的核验清单；尤其补齐执行时的锁定期、卡状态等业务重验，明确各入口的限额、版本和时限，不依赖规划时曾经检查过的条件。
\item \textbf{细化任务终态与恢复。} 将整单执行尝试结束与全部成功区分，面向用户展示部分完成、失败及暂停项；重新办理时引用原回执，避免补做任务导致重复。
\item \textbf{收紧对外数据控制。} 对当前消息、历史、目标名称和检索词提供清楚的发送提示与可测试的字段治理；对真实美团地址、供应商响应和缓存建立保留与访问策略。
\item \textbf{分开演示与真实授权空间。} 共享 persona 用于合成演示，真实外部账号使用独立注册身份。公网配置平台密钥应采用明确管理员身份与最小权限，而不是仅依靠本机开关。
\item \textbf{完善目标环境验证。} 固化 Linux 镜像与官方组件版本，补 Windows、真实手机及外部服务验收，建立备份恢复和关键状态观测。SQLite 单机适用于当前演示；多实例演进需重新设计数据库锁、共享任务与幂等边界。
\end{enumerate}

\subsection{面向更完整金融生活的方向}
项目可以继续探索家庭协同、跨账户资金视图、小微经营中的应收应付与回款压力，以及更长时期的目标衔接。前提是获得相应数据与授权，明确金融业务资质、供应商契约和可验证的决策边界。不能把当前 30 天有限模型外推为已经实现的完整生命周期优化。

TROPOS 的长期方向仍是同一件事：理解人的处境，让安排跟随变化，同时把决定权与可追溯的结果留给用户。技术演进应以这一要求衡量，而不是以增加模型数量或场景数量作为完成标准。
